% Clinical Background chapter. Paste into your own thesis.
%
% Needs only graphicx (for the one figure) -- no siunitx, no cleveref, no booktabs.
% Citation keys are defined in refs.bib in this folder; append those entries to your own .bib.
% The figure is ImageCAS-X/figures/05_dominance_comparison.png -- adjust the path or delete
% the figure block if you do not want it.

\chapter{Clinical Background}
\label{ch:clinical}

This chapter establishes what the coronary arteries are, how they vary between people, how
computed tomography records them, and why the shape of the arterial tree carries clinical
meaning.
Each section is included because a later part of this thesis depends on it.
Sections~\ref{sec:heart} and~\ref{sec:tree} fix the vocabulary that the segmentation labels
use.
Sections~\ref{sec:dominance} and~\ref{sec:variants} describe the variation that the
topological characterization must accommodate rather than treat as error.
Sections~\ref{sec:cad} and~\ref{sec:ccta} give the disease and the instrument.
Section~\ref{sec:geometry} states the mechanism by which a geometric measurement becomes a
clinical one, and is what licenses the feature extraction developed in later chapters.

Where a quantity in this chapter has also been measured in the cohort used by this thesis,
both values are given.
The cohort numbers were computed over the 800 usable cases of ImageCAS-X
\cite{bransby2026imagecasx} and are reproducible with the scripts named in each case.

\section{The heart and the coronary circulation}
\label{sec:heart}

The heart is a four-chambered muscular pump that drives two circuits in series
\cite{weinhaus2024anatomy}.
Deoxygenated blood returns from the body through the venae cavae into the right atrium,
passes into the right ventricle, and is pumped through the pulmonary circuit where it is
oxygenated.
It returns to the left atrium, passes into the left ventricle, and is ejected through the
aorta to the rest of the body.
Four valves enforce the direction of flow: the tricuspid and mitral valves between atria and
ventricles, and the pulmonary and aortic valves between ventricles and arteries
\cite{weinhaus2024anatomy,shah2009heart}.

The cardiac cycle alternates between diastole, when a chamber relaxes and fills, and systole,
when it contracts and ejects \cite{shah2009heart}.
Atria and ventricles alternate, so that ventricular diastole coincides with atrial systole.
This alternation matters for imaging rather than only for physiology, because the coronary
arteries move with the myocardium and are imaged during the quiescent phase of the cycle
(Section~\ref{sec:ccta}).

The muscular wall that performs the contraction is the myocardium, the thickest of the three
layers of the heart wall \cite{weinhaus2024anatomy}.
The myocardium is not supplied by the blood passing through the chambers.
It has its own arterial supply, the coronary circulation, and the consequence of that
separation is the subject of this thesis: an obstruction in a vessel a few millimeters across
can arrest a pump that is otherwise intact.

\section{The coronary arterial tree}
\label{sec:tree}

The coronary arteries arise from the aorta immediately above the aortic valve, from two of
the three dilatations of the aortic root known as the aortic sinuses
\cite{devos2024coronary}.
The left coronary ostium opens from the left sinus and the right coronary ostium from the
right sinus.
Each ostium is the root of an arterial tree, and the two trees are treated separately
throughout this thesis because the dataset stores them as separate structures
\cite{bransby2026imagecasx}.

\subsection{The left system}

The left coronary artery begins as the left main artery (LM), a short trunk that typically
divides into the left anterior descending artery (LAD) and the left circumflex artery (LCX)
\cite{devos2024coronary}.
The LAD runs in the anterior interventricular groove toward the apex and supplies the
anterior wall of both ventricles and the anterior two-thirds of the interventricular septum,
by way of septal perforators and a variable number of diagonal branches
\cite{devos2024coronary,obrien2007anatomy}.
The LCX runs in the left atrioventricular groove and supplies the lateral wall of the left
ventricle through one or more obtuse marginal branches \cite{devos2024coronary}.

The diagonal and obtuse marginal branches are numbered by their order of origin along the
parent vessel rather than by any fixed identity: D1 is simply the first diagonal arising from
the LAD.
This convention has a measurable consequence in the cohort, where prevalence decays with
ordinal position.
D1 is present in 98.5\% of cases and D2 in 59.4\%; OM1 is present in 86.8\% and OM2 in
37.6\%.

\subsection{The right system}

The right coronary artery (RCA) runs in the right atrioventricular groove
\cite{devos2024coronary}.
Along its course it gives off the conus and sinoatrial nodal branches proximally and the
acute marginal branch to the right ventricular free wall, before terminating in the posterior
descending artery (PDA) and one or more posterior lateral arteries (PLA) in most people
\cite{devos2024coronary,obrien2007anatomy}.
The PDA supplies the inferior third of the interventricular septum.

The proximal branches of the RCA are not represented in the ImageCAS-X label set, which names
only the RCA trunk and the posterior vessels on the right side.
This is an annotation boundary rather than an anatomical claim, and it is stated here because
it constrains any comparison of branch counts between the left and right trees later in this
thesis.

\subsection{Territories}

The clinical weight of a lesion follows from the territory distal to it rather than from the
lesion itself.
An occlusion of the proximal LAD jeopardizes the anterior wall and the septum, a large
fraction of the left ventricular mass, whereas an occlusion of a distal obtuse marginal
jeopardizes a small lateral segment \cite{devos2024coronary}.
The same reasoning gives distance from the ostium its status as a risk-relevant geometric
quantity, a point returned to in Section~\ref{sec:geometry}.

\section{Coronary dominance}
\label{sec:dominance}

Coronary dominance is defined by the origin of the posterior descending artery
\cite{allwork1987applied}.
A heart is right-dominant when the PDA arises from the RCA, left-dominant when it arises from
the LCX, and codominant when the posterior territory is shared between the two systems.
Dominance is the single largest source of normal topological variation in the coronary tree,
which is why it is treated as a stratifying variable throughout this thesis rather than as a
nuisance parameter.

Reported prevalences vary widely: right dominance between 70\% and 89\%, left dominance
between 8\% and 12\%, and codominance anywhere from 2.5\% to 17.5\%
\cite{jaha2024dominance}.
The spread in the codominance figure is not measurement error.
It reflects genuinely different definitions of what constitutes a shared posterior supply
across studies \cite{jaha2024dominance}, which means any cohort figure is only interpretable
alongside the definition used to produce it.

The cohort used here is annotated as 729 right-dominant, 41 left-dominant and 30 codominant
cases, or 91.1\%, 5.1\% and 3.8\% of the 800 usable cases \cite{bransby2026imagecasx}.
Right dominance is therefore more common here than in the ranges above, and codominance
rarer.

Dominance is recorded twice in this dataset, independently, and the two records agree.
It appears as an explicit annotation, and it is implied by which arteries are labeled, since
the posterior vessels arise from whichever system is dominant.
Table~\ref{tab:posterior} gives the cross-tabulation.
The right posterior vessels appear in essentially every right-dominant case and in no
left-dominant case; the left posterior vessels invert the pattern exactly.
Codominant cases draw on both, which is the definitional ambiguity above made visible in the
data.

\begin{table}[t]
  \centering
  \caption{Prevalence of the posterior vessels by annotated dominance, over the 800 usable
    cases. The right and left posterior vessels are mutually exclusive in the
    single-dominance groups, so the artery labels and the dominance annotation are two
    independent records of the same anatomical fact.}
  \label{tab:posterior}
  \begin{tabular}{lrrrr}
    \hline
    Vessel & All (\%) & Right (\%) & Left (\%) & Co (\%) \\
    \hline
    R-PDA & 93.0 & 99.7 &  0.0 & 56.7 \\
    R-PLA & 92.0 & 99.2 &  0.0 & 43.3 \\
    L-PDA &  5.8 &  0.0 & 90.2 & 30.0 \\
    L-PLA &  8.8 &  1.0 & 90.2 & 86.7 \\
    \hline
  \end{tabular}
\end{table}

\begin{figure}[t]
  \centering
  \includegraphics[width=\textwidth]{05_dominance_comparison.png}
  \caption{Right-, left- and codominant coronary trees from the cohort. The posterior
    vessels follow the dominant system, which is why the right posterior descending artery is
    absent from the left-dominant tree entirely rather than merely smaller.}
  \label{fig:dominance}
\end{figure}

\section{Anatomical variants}
\label{sec:variants}

Beyond dominance, several variants change the connectivity of the coronary tree rather than
only its dimensions.
They are described here because a method that characterizes coronary topology must represent
them rather than discard them, and because two of them occur often enough in this cohort to
affect population statistics.

\subsection{Absent left main}

In a minority of people the left main artery does not exist as a discrete trunk, and the LAD
and LCX arise instead from separate ostia in the left aortic sinus
\cite{devos2024coronary,obrien2007anatomy}.
The left tree then has two roots rather than one.

The variant is present in 11 of the 800 cases, or 1.4\%, and it is identifiable from two
independent signals that agree exactly: those cases carry no LM label, and their left
centerlines resolve into two components each with its own ostium.
The practical consequence is that no stage of the analysis may assume a single ostium per
side.

\subsection{Ramus intermedius}

When the left main trifurcates rather than bifurcates, the additional vessel arising between
the LAD and the LCX is the ramus intermedius, labeled IM in this dataset
\cite{obrien2007anatomy}.
It supplies territory that would otherwise belong to a diagonal or an obtuse marginal branch.

The reported prevalence and the cohort prevalence disagree.
O'Brien et al.\ report trifurcation of the left main in approximately 15\% of people
\cite{obrien2007anatomy}, whereas an IM label is present in 196 of the 800 cases here, or
24.5\%.
The cohort figure is stable across dominance groups at 24.7\%, 22.0\% and 23.3\% for right-,
left- and codominant cases respectively, which makes a sampling artifact within the cohort
unlikely.
The discrepancy is therefore attributable either to the source population or to the threshold
at which an annotator records a third branch as a ramus intermedius rather than a proximal
diagonal.
It is left open here and revisited when the population statistics are presented.

\subsection{Myocardial bridging and anomalous origin}

A segment of a coronary artery may run within the myocardium rather than on its surface, a
variant termed myocardial bridging, most often affecting the mid-LAD
\cite{obrien2007anatomy}.
The bridged segment is compressed during systole.
Separately, a coronary artery may arise from the wrong sinus or from the pulmonary artery, a
group of anomalies whose clinical significance depends on whether the anomalous vessel takes
a course between the aorta and the pulmonary trunk \cite{obrien2007anatomy}.
Neither variant is separately labeled in this dataset, so neither can be studied directly
here; both are noted because they are candidate explanations for individual outliers in later
chapters.

\section{Coronary artery disease}
\label{sec:cad}

Cardiovascular disease remains the leading cause of death worldwide, accounting for
approximately 20.5 million deaths in 2021 \cite{lindstrom2022global}.
Coronary artery disease (CAD) is the largest component of that total.

CAD is the manifestation of atherosclerosis in the coronary arteries.
Atherosclerosis begins with injury to the endothelium, the single cell layer lining the
arterial lumen, followed by accumulation of lipid, inflammatory cells and eventually calcium
within the vessel wall \cite{geometry2022plaqueburden}.
The resulting plaque may narrow the lumen and restrict flow.

Two distinctions matter for what follows.
The first is between plaque composition types, since calcified, non-calcified and mixed
plaques differ in their risk of rupture, and the non-calcified plaques that are hardest to
see are the more dangerous \cite{geometry2022plaqueburden}.
The second is between anatomical stenosis and functional ischemia: a narrowing that appears
severe on an image does not necessarily reduce flow enough to starve the myocardium, which is
why the diagnostic pathway distinguishes anatomical from functional testing
\cite{knuuti2019esc}.

The acute event that ends this process is usually not the gradual closure of a lumen but the
rupture of a plaque, exposing thrombogenic material and forming a clot that occludes the
vessel abruptly \cite{geometry2022plaqueburden}.
The myocardium distal to the occlusion is deprived of oxygen and infarcts.
This is why the location and character of plaque, rather than the degree of narrowing alone,
predict clinical events, and it is the opening through which vessel geometry becomes
prognostically relevant (Section~\ref{sec:geometry}).

\section{Coronary computed tomography angiography}
\label{sec:ccta}

Coronary computed tomography angiography (CCTA) is the non-invasive imaging test recommended
as an initial investigation in patients with suspected chronic coronary syndrome
\cite{knuuti2019esc}.
An iodinated contrast agent is injected and the scan is timed so that acquisition coincides
with peak opacification of the coronary lumen.
The result is a three-dimensional volume in which the contrast-filled lumen is bright against
surrounding tissue.

This section treats CCTA as a measuring instrument, because the limits of the instrument
bound every geometric claim made later in this thesis.

\subsection{Acquisition and its failure modes}

The coronary arteries move throughout the cardiac cycle, so acquisition is synchronized to
the electrocardiogram and data are reconstructed at the quiescent phase, conventionally in
mid-diastole.
Heart rate is commonly lowered pharmacologically before the scan to lengthen that quiescent
window.
When synchronization fails, or when the rhythm is irregular, motion artifact blurs or
duplicates the vessel.

Dense calcium produces blooming, in which the apparent extent of a calcified plaque exceeds
its true extent and the adjacent lumen is obscured.
Blooming is the reason a heavily calcified segment may be uninterpretable even in an
otherwise excellent scan.

\subsection{Resolution against vessel caliber}

The volumes used in this thesis have in-plane spacing between 0.318 mm and 0.43 mm and a
slice thickness of 0.5 mm, varying between cases \cite{zeng2023imagecas}.
Proximal coronary arteries are a few millimeters in diameter and taper distally toward 1 mm
and below.
A vessel two or three voxels across cannot support a confident radius measurement, and this
places a floor on the branch order at which geometric features remain meaningful.

Contrast opacification in this cohort is adequate: sampling the reconstructed volumes along
the annotated centerlines gives a median of 336 HU, consistent with a contrast-filled lumen.

\subsection{Image quality as a variable}

Image quality is graded from 0 to 4 in this dataset, and the 200 cases graded 0 are excluded
from the released cohort, leaving the 800 used here \cite{bransby2026imagecasx}.
Among those 800 the grading remains informative rather than uniform.
Because a lower-quality scan plausibly yields fewer traceable distal branches, image quality
is a candidate confounder for any branch-count statistic and is retained as a stratifying
variable.

\section{Geometry as a risk factor}
\label{sec:geometry}

The preceding sections describe the coronary tree and the disease that affects it.
This section states why the \emph{shape} of the tree is a clinical quantity rather than an
incidental descriptor, and it is the argument on which the remainder of this thesis rests.

The accepted causal chain runs from geometry through hemodynamics to biology: geometry and
topology determine wall shear stress, wall shear stress drives the endothelial response, and
the endothelial response produces plaque.

Wall shear stress (WSS) is the tangential frictional force exerted by flowing blood on the
endothelium.
Regions of low and oscillatory shear are where early atherosclerotic lesions preferentially
develop, because the endothelium responds to disturbed shear with a pro-inflammatory
phenotype \cite{geometry2022plaqueburden}.
The geometric configurations that produce low shear are predictable.
Flow separates at a bifurcation, leaving a low-shear region on the lateral walls away from
the flow divider, and curvature drives secondary flows that leave a low-shear region on the
inner wall of a bend \cite{geometry2022plaqueburden}.
Bifurcation angle and tortuosity are therefore not arbitrary descriptors of shape; they are
proxies for the size and location of the regions in which lesions form.

This reasoning is supported by outcome data.
In a study of patients who subsequently experienced an acute coronary syndrome, three
geometric characteristics carried prognostic value beyond stenosis severity and plaque
characteristics: a more proximal lesion location, location at a bifurcation, and increased
tortuosity \cite{han2022plaque}.
That result is the strongest available justification for treating the geometry of the
coronary tree as an object of study in its own right.

Two quantitative anchors are worth recording, because later chapters compare against them.
The normal left main bifurcation angle has been reported as $76.4^\circ \pm 16.7^\circ$, and
curvature has been reported as approximately 16.7\% higher in stenotic than in non-stenotic
segments \cite{geometry2022plaqueburden}.
The direction of the association between bifurcation angle and risk is not consistent across
the literature, so no sign is assumed in this thesis; the measured association is reported as
found.

A scaling relation completes the picture.
Murray's law, derived by minimizing the sum of the power dissipated in flow and the metabolic
cost of maintaining blood, predicts that flow scales with the cube of vessel radius, and
hence that wall shear stress is constant throughout a healthy arterial tree.
In the coronary circulation specifically the exponent is not three.
A meta-analysis over 1070 coronary trees places it at 2.39, with a 95\% confidence interval
of 2.24 to 2.54 \cite{taylor2024murray}.
Any feature defined as a deviation from Murray's law must use the coronary exponent, and this
thesis does so.

Finally, the approach taken here has a direct antecedent.
Wang et al.\ extracted centerline geometry from 76 left main arteries, clustered the cases
into four morphological phenotypes, and computed wall shear stress for each phenotype,
finding reduced shear near the LAD branch point in the cluster characterized by a large
bifurcation angle \cite{wang2024lmphenotypes}.
The present work differs in scale, in extent and in framing: 800 cases rather than 76, the
whole coronary tree rather than the left main alone, and the detection of population-level
outliers rather than the recovery of cluster centroids.
